\cotefile{156-6}
\scanpage{88}\mspage{87}

\uline{At \deleted{\unsure{14 bis}}(\unsure{fac|Fac}) 14 bis} Sous les conditions précédentes, \deleted{\unsure{considérons}} considérons
\deleted{\unsure{Mettons le|tel que le|tels que le}}
\diagram{\begin{tikzcd}[ampersand replacement=\&, column sep=.3em,row sep=1.3em]
\widetilde G \& \subset \& \widetilde F \\
\widetilde G' \arrow[u] \& \&
\end{tikzcd}}
alors on peut trouver $F'\preccurlyeq F$ induisant $G'$, de telle façon que dans le diagramme d'ens \inserted{ordonnés}
\diagram{\begin{tikzcd}[ampersand replacement=\&,column sep=2.6em,row sep=2.6em]
\widetilde G'_{\deleted{\unsure{t|+}}} \arrow[r,hook] \arrow[d,"{\varphi\mathbin{\unsure{\vert}}\widetilde G'}"'] \arrow[dr,phantom,"\text{\unsure{cart}}" description] \& \widetilde F' \arrow[d,"\varphi"]\\
\underset{\substack{\text{\unsure{sous ens.}}\\\text{fermé}}}{\widetilde G} \arrow[r,hook] \& \widetilde F
\end{tikzcd}}
on ait $\widetilde F'\setminus\widetilde G'_{t}\overset{\substack{\text{iso d'ens.}\\\text{ordonnés}}}{\xrightarrow{\sim}}\widetilde F\setminus\widetilde G$, et pour $\xi\in\widetilde G'_{t}$, $\eta\in\widetilde F'$, on ait $\xi\leq\eta$ \unsure{sss|ssi} $\varphi^{\unsure{\#|*}}(\unsure{\vert|\xi})\leq\varphi^{\#}(\eta)$.

[ce qui détermine l'ens. ordonné $\widetilde F'$ et le diagramme précé\uline{dent,} à iso unique près].

\margin{\uline{22. juin}}Je présume qu'on doit pouvoir prouver qu'un tel $F'$ est unique, comme la plus grande des subdivisions de $F$ induisant $G'$ sur $G$.\margin{\unsure{unicité automatique} ???}% slanted
 Dans \unsure{les|le} cas où $\unsure{\mathcal F|\exists}$ est à figures finies (pour \unsure{éviter|iviter} des inductions transfinies \ldots), \deleted{cette} condition doit être équivalente à

\noindent\begin{tabular}{@{}||p{.90\linewidth}@{}}
At (fac) 14 ter : Soit $X\in M$, et $F'\preccurlyeq\partial X$. Alors il existe un unique $X'\in M$ tel que $X'\preccurlyeq X$, $\partial X'=F'$
\end{tabular}
